% =============================================================================
% Project : THE_PHILOSOPHY_OF_VIABILITY
% Book    : Book I — Ontology of Viability
% Book ID : VS.B01
%
% File    : 4_ASSETS/tables/VS.B01.P01.C03.S07.SS10.P003.TL01.unity_singularity_and_related_concepts.tex
%
% Asset:
%   Table TL01 — Unity, Singularity, and Related Concepts
%
% Purpose :
%   To distinguish ontological unity from numerical singularity,
%   uniqueness, indivisibility, homogeneity, individuality,
%   individuation, and the number of identities.
%
% =============================================================================

\begin{table}[htbp]
\centering
\caption{Unity, Singularity, and Related Concepts}
\label{tab:unity-singularity-and-related-concepts}
\footnotesize
\renewcommand{\arraystretch}{1.2}
\begin{tabularx}{\textwidth}{
    @{}
    p{0.18\textwidth}
    X
    p{0.25\textwidth}
    p{0.27\textwidth}
    @{}
}
\toprule
\textbf{Concept} &
\textbf{Principal Question} &
\textbf{Required Specification} &
\textbf{What It Does Not Establish}
\\
\midrule

Numerical singularity &
How many units of a specified kind are counted? &
Counting domain, unit-kind, and criterion &
Ontological unity or the number of identities
\\

Ontological unity &
In what respect is a subject unified? &
Subject, respect, scope, and relevant grounds &
Numerical singularity or exactly one identity
\\

Uniqueness &
Is there another item satisfying the same criterion? &
Comparison domain and distinguishing criterion &
Unity, indivisibility, or a unique identity
\\

Indivisibility &
Can the subject be divided under a specified criterion? &
Kind and level of division &
Unity of identity or universal indivisibility
\\

Homogeneity &
Are elements similar in a specified respect? &
Shared feature and respect of comparison &
Unity of the subject or identity
\\

Individuality and individuation &
What makes a subject an individual or individuated? &
Relevant principle or criterion of individuation &
Unity, numerical singularity, or identity as such
\\

Unity of identity &
What ontological unity is attributed to identity? &
Meaning, subject, respect, and scope of the unity claim &
A numerical conclusion unless that conclusion is separately established
\\

Number of identities &
How many identities are involved? &
Identity-related grounds of distinction and relevant domain &
A result inferred solely from component, manifestation, or aggregate counts
\\

\bottomrule
\end{tabularx}
\end{table}

% =============================================================================
% END OF TABLE TL01
% =============================================================================